% SedSAT3 – Response to the domain review (L. Gorman Sanisaca, October 2026)
% Build: pdflatex SedSAT3_domain_review_response.tex  (run twice)
% Word copy: python3 ../response-letter/tex2docx.py SedSAT3_domain_review_response.tex SedSAT3_domain_review_response.docx
\documentclass[11pt]{article}
\usepackage[letterpaper,margin=1in]{geometry}
\usepackage[T1]{fontenc}
\usepackage{lmodern}
\usepackage{microtype}
\usepackage[utf8]{inputenc}
\usepackage{parskip}
\usepackage{xcolor}
\usepackage{enumitem}
\usepackage{booktabs}
\usepackage{longtable}
\usepackage{tabularx}
\usepackage[most]{tcolorbox}
\usepackage{fancyhdr}
\usepackage[hidelinks]{hyperref}

\definecolor{usgsgreen}{RGB}{0,100,60}
\definecolor{commentbg}{RGB}{242,245,248}
\definecolor{todo}{RGB}{190,30,30}
\definecolor{mine}{RGB}{20,70,170}

% ---------------------------------------------------------------- macros
\newcommand{\code}[1]{\texttt{#1}}
\newcommand{\commit}[1]{\texttt{#1}}
\newcommand{\todo}[1]{\textcolor{todo}{\textbf{[#1]}}}   % for Cara / Allen
\newcommand{\mine}[1]{\textcolor{mine}{\textbf{[Arash: #1]}}}   % for the author
\newcounter{cmt}
\newcounter{grp}
\renewcommand{\thecmt}{\thegrp.\arabic{cmt}}
\newcommand{\group}[1]{\stepcounter{grp}\setcounter{cmt}{0}\section*{\thegrp.\quad #1}\addcontentsline{toc}{section}{\thegrp.\quad #1}}
\newenvironment{reviewer}[1]{%
  \refstepcounter{cmt}%
  \begin{tcolorbox}[enhanced,breakable,colback=commentbg,colframe=usgsgreen!60,
    boxrule=0.4pt,left=6pt,right=6pt,top=4pt,bottom=4pt,
    title={\small\textbf{Comment \thecmt}\hfill\textit{#1}},
    coltitle=black,colbacktitle=usgsgreen!15,fonttitle=\small]}%
  {\end{tcolorbox}}
\newenvironment{response}{\par\noindent\textbf{Response.}\ }{\par\medskip}
\newcommand{\status}[1]{\par\noindent\textit{Status:} #1\par}

\pagestyle{fancy}
\fancyhf{}
\lhead{\small SedSAT3 — response to the domain review}
\rhead{\small \thepage}
\renewcommand{\headrulewidth}{0.3pt}

\begin{document}
\sloppy
\setlength{\emergencystretch}{3em}

\begin{tcolorbox}[colback=todo!4,colframe=todo!70,boxrule=0.6pt]
\textbf{Draft.} Items in \textcolor{todo}{\textbf{red}} need a decision from
Cara Peterman-Phipps or Allen Gellis; items in \textcolor{mine}{\textbf{blue}}
are tasks for Arash. All marked items and this note will be removed before the
response is final.
\end{tcolorbox}

\begin{flushleft}
Arash Massoudieh\\
Department of Civil and Environmental Engineering\\
The Catholic University of America\\
\href{mailto:massoudieh@cua.edu}{massoudieh@cua.edu}

\medskip
\today

\medskip
Lillian Gorman Sanisaca\\
U.S. Geological Survey
\end{flushleft}

\textbf{Re: Domain review of SedSAT3 (Sediment Source Assessment Tool, version 3)}

Dear Lilly,

Thank you for your careful domain review of SedSAT3: the review notes and the
125 comments in the user's manual. You tested every tool, built test cases for
the bracket test, the column layout and the column names, and recorded each
crash with the settings that caused it. That made every problem reproducible.
This letter answers each of your notes and comments in turn.

\textbf{The version you tested.} Your copy was the Windows package of release
1.1.6, built in December 2025. Since then, the code review by Steve Westenbroek
(October 2026) led to a large set of fixes on the branch
\code{review-fixes-oct2026} (merge request !4 on code.usgs.gov). Several of the
problems you found were already fixed there: copies of results that lost their
red highlighting and axis titles, progress windows that stayed open, a chart
crash on negative values, and distribution fitting. I reproduced every one of
your crashes against the current code with the example data, and fixed those
that remained. The changes made for your review are in four commits on the
same branch: \commit{d7d8b92}, \commit{8708c20}, \commit{ae2e4d5} and the
manual commits.

\textbf{What has changed, in brief.}
\begin{itemize}
    \item \textbf{All seven crashes are fixed} (Section~1). Most were caused by
          one problem: after step~4 with \textit{Modify the included elements},
          the base elements of the isotopes (Calcium and N) were excluded while
          the isotopes stayed. The fingerprinting tools now stop with a message
          that names the isotope and its base element. Element selection no
          longer excludes the particle size and organic matter constituents.
    \item \textbf{D50 and TOC are no longer treated as tracers.} The tools that
          test or select tracers never include particle size or organic matter
          constituents.
    \item \textbf{The Tools menu} follows the workflow order, and Other
          Statistical Tools is grouped: Data exploration, Element selection,
          Discriminant function analysis, Uncertainty and verification.
    \item \textbf{Charts}: a colour-blind-safe palette with the same colour for
          a group in every plot of a window; axis titles that name what is
          plotted; axis ranges that include every point; p-value plots that
          always show their axis; observed values in black.
    \item \textbf{Step~4} also shows the discriminant scores and eigenvectors of
          the constituents it selects.
    \item \textbf{Bayesian defaults}: 100,000 samples and 20,000 burn-in.
    \item \textbf{The user's manual} has been revised for every comment
          (Sections~5 to~9). Its figures were regenerated from the example data
          by the current program, at full page width. Each tool's description in
          the program ends with the manual section that documents it.
\end{itemize}

The revised manual is attached as a PDF (\code{SedSAT3\_User\_Manual.pdf}); its
LaTeX source is in \code{docs/manual/SedSat\_User's\_Manual} in the repository.
Section numbers below refer to it. A few items need a decision from Cara and
Allen; they are listed in Section~10.

With thanks,

Arash Massoudieh

\newpage
\tableofcontents
\newpage

% =====================================================================
\group{Critical errors (program closed)}

For each crash I recreated your dataset and settings (the example data with
Aluminum in column 2, the renamed columns, CTAIL1/Aluminum set to 10000, Lead
excluded, the size and organic matter correction for CTAIL1) and ran the tool
with the current code.

\begin{reviewer}{Review notes}
Tools → Other Statistical Tools → Correct based on size and organic matter,
graph clicked.
\end{reviewer}
\begin{response}
Already fixed before your review, in commit \commit{acfc109}. With the linear
correction, some corrected values (Manganese, Cobalt) are negative. The chart
then fell back to a linear axis but attached the points to a logarithmic axis
that it had never created. The chart now uses the axis it creates. The program
also reports the negative corrected values (Comment~7.4).
\end{response}
\status{Fixed.}

\begin{reviewer}{Review notes}
Deterministic Sediment Fingerprinting → 5-Fingerprinting and 6-Fingerprinting
(batch).
\end{reviewer}
\begin{response}
Reproduced with the current code. As you found, the cause was that the base
elements of the isotopes (Calcium for C13 and N for N15) had been excluded,
here by step~4 with \textit{Modify the included elements}. Fingerprinting
needs the base element's concentration to turn a $\delta$ value into an
isotope concentration, and it read a distribution that did not exist.
Fingerprinting, its batch version, the Bayesian tools, Error Analysis and
Source Verification now check every included isotope first. When a base
element is excluded, they stop with a message: ``The base element `N' of
isotope `N15' is excluded from the analysis. Include the base element in the
analysis (Settings > Constituent Properties), or exclude the isotope.'' The
manual says so where isotopes are set up (Section~2.1) and in Section~3.6.
Unit test \code{fingerprinting\_isotopeBaseElementExcluded}.
\end{response}
\status{Fixed (\commit{d7d8b92}).}

\begin{reviewer}{Review notes}
Other Statistical Tools → Distribution Fitting.
\end{reviewer}
\begin{response}
Already fixed in the code review work (commit \commit{ac039a2}). An internal
series was created under the wrong name, and the lookup of ``Observed'' then
failed. Your four settings (Box--Cox on and off; C13, Be and Ba; selected
samples on and off; the correction for CTAIL1) all run now. For C13 without
Box--Cox, the tool stops with a message, because a log-normal distribution
cannot be fitted to negative values.
\end{response}
\status{Fixed.}

\begin{reviewer}{Review notes}
Other Statistical Tools → Elements' Discriminant Power (Multi-way).
\end{reviewer}
\begin{response}
Reproduced with the current code. The multi-way tool built its title from two
source groups, as the two-way tool does, but its form has no group fields, so
it failed on every setting. Fixed. Unit test
\code{multiwayDiscriminantPower\_runs}.
\end{response}
\status{Fixed (\commit{d7d8b92}).}

\begin{reviewer}{Review notes}
Other Statistical Tools → Error Analysis; Other Statistical Tools → Source
Verification.
\end{reviewer}
\begin{response}
Both run with the current code when the isotopes' base elements are included;
they failed in your session for the reason given above (Calcium and N
excluded). They now stop with the same message instead. While checking Error
Analysis I also found that its frequency-distribution panel plotted only the
first realization, a narrow spike. It now uses all realizations
(\commit{8708c20}).
\end{response}
\status{Fixed.}

\begin{reviewer}{Comment in the manual, Elements' correlation matrix}
Attempted to run with OM and Size correction based on CTAIL1, all checkboxes
clicked: program closed. Works with no corrections.
\end{reviewer}
\begin{response}
Reproduced with the current code. With \textit{Use only selected elements},
the tool dropped the excluded constituents, including TOC and D50, before
applying the correction that needs them. It now corrects first and filters
afterwards, as the other tools do. Unit test
\code{correlationMatrix\_correctedWithSelectedElements}.

The test for this fix also exposed another crash: the organic matter and size
regression stopped the program (inside the GSL library) when a group had no
more samples than regression coefficients. Such a group now gets a regression
that corrects nothing.
\end{response}
\status{Fixed (\commit{d7d8b92}).}

% =====================================================================
\group{README}

\begin{reviewer}{Review notes}
Neither the readme nor the documentation have up-to-date instructions on how to
install and run the program. Instructions are only available in the readme
within the zip file. Both the front page readme on code.usgs.gov and the
documentation should have basic setup instructions.
\end{reviewer}
\begin{response}
The README on the review branch, which becomes the front page when merge
request !4 is merged, has a section ``Running the portable Windows build'' and
build instructions for every platform. The manual has a new chapter, ``Getting
started'' (Section~1.1), with the same installation steps.
\end{response}
\status{Done.}

\begin{reviewer}{Review notes}
Opening paragraph ends in ``\ldots'' — is it incomplete? Monte Carlo should be
spelled out, and (MCMC) put after the full name.
\end{reviewer}
\begin{response}
The README was rewritten during the code review. Its opening paragraph is
complete, and it spells out ``Markov chain Monte Carlo (MCMC)'' at first use.
\end{response}
\status{Done.}

\begin{reviewer}{Review notes}
Links to GitHub will need to be updated to code.usgs.gov. Citation needs a DOI,
not the GitHub URL \url{https://sedsat.org/}.
\end{reviewer}
\begin{response}
The README no longer links to GitHub or sedsat.org. Its repository link and
citation hold placeholders for the final code.usgs.gov address and the
software release DOI, which are set at release.
\todo{Cara: the final repository path and the DOI.}
\end{response}
\status{Placeholders in place; values pending the release.}

\begin{reviewer}{Review notes}
I do not know the regulations around storing software on a non-government
public website, or around online forms that collect name, email and company.
Please verify that this meets software release and online form guidelines.
Is \url{https://sedsat.org/} just for the Windows download? Its MacOS and Linux
links: do they bypass the Readme instructions? The website gave a warning that
the download may not be safe.
\end{reviewer}
\begin{response}
These questions are about sedsat.org, a site I set up before the USGS release.
The USGS release will be distributed from code.usgs.gov, and the README and the
manual no longer refer to sedsat.org.
\todo{Cara, Allen: decide whether sedsat.org is retired, or kept only as a
pointer to the USGS release, and check its download form against USGS
guidelines if it stays.} The browser warning is likely because the Windows
executable is not code-signed; packages from code.usgs.gov have the same
issue unless they are signed. \todo{Cara: ask the release coordinator whether
the Windows package should be signed.}
\end{response}
\status{Decision for Cara and Allen.}

% =====================================================================
\group{Sample data}

\begin{reviewer}{Review notes}
I found three example data files. The two called SampleData.xlsx were
identical; virtual-mixture.xlsx (in the package registry) differs: Calcium is
``Calcium\_µg/g'', and Nickel, Potassium, Uranium and Vanadium have ``gg''
appended.
\end{reviewer}
\begin{response}
The example dataset for the manual and the release is \code{examples/SampleData.xlsx}
in the repository. The manual now names it and no longer links to sedsat.org
(Section~1.1). \code{virtual-mixture.xlsx} was a test file with synthetic target
samples. Its odd column names come from an old test of unusual characters in
names, and it should not be shipped as an example.
\mine{remove \code{virtual-mixture.xlsx} from the Windows package and include
\code{examples/SampleData.xlsx} instead.}
\end{response}
\status{Manual done; package to update.}

% =====================================================================
\group{The program}

\begin{reviewer}{Review notes}
Many steps have the checkboxes ``Use only selected elements'' and ``Use only
selected samples'', but nothing stops the user from selecting yes in one part of
the analysis and no in another. Are the user's selections saved somewhere?
\end{reviewer}
\begin{response}
These options are intentionally per tool: an exploratory tool can look at all
constituents while fingerprinting uses only the selected ones. The manual now
explains this in a new section, ``Settings shared by many tools''
(Section~1.3.2). It says that each tool applies the options to its own run only,
that results compare only when run with the same options, and that the
selections are saved in the project file. Two changes make the selections
safer. Particle size and organic matter constituents can no longer be selected
or excluded as tracers by any tool. The fingerprinting tools refuse an isotope
whose base element is excluded.
\end{response}
\status{Documented; partly changed.}

\begin{reviewer}{Review notes}
After an error in DFA with linear corrections, I went back to step 2 and chose
power transforms, but at the bracket test step only 6 tracers appeared.
Warnings should be given in the program and the documentation that returning to
a previous step can result in mixups.
\end{reviewer}
\begin{response}
The six tracers were those selected by step~4: \textit{Modify the included
elements} had excluded the others, and that exclusion applies to every later
tool, including earlier steps run again. The warning is now in the tool
descriptions of step~4, ANOVA and auto-select, and in a note in Section~1.3.2
of the manual.
\end{response}
\status{Done.}

\begin{reviewer}{Review notes}
In many cases, once a step has completed, there is a progress window at 100\%.
Could it close automatically when the results window is closed? (Also: the
progress window appears behind the results window; users may think the program
failed.)
\end{reviewer}
\begin{response}
Already fixed in the code review work. Progress windows now close as soon as
the analysis ends, so none is left behind the results.
\end{response}
\status{Fixed.}

\begin{reviewer}{Review notes}
Several plots have points at the axis limits; expand the axes. Check axis labels
and tick marks on all plots.
\end{reviewer}
\begin{response}
I went through every chart type, rendering each with the example data. The
changes:
\begin{itemize}
    \item The axes now extend beyond the data: the lower limit was rounded to
          the nearest value instead of down, which cut off the smallest points
          and bars.
    \item The axis titles name what is plotted: ``Constituent'',
          ``Constituents, in the order added'', ``Discriminant score'',
          ``Probability density'', and the name of the quantity rather than
          ``Value''. Several plots had none.
    \item MCMC sample numbers are whole numbers.
    \item The Box--Cox plot has a grid line through zero.
    \item The p-value plots always show their axis. Qt's logarithmic axis
          disappears for values below about $10^{-12}$, so such values are now
          plotted as base-10 logarithms, which the axis title says.
\end{itemize}
\end{response}
\status{Fixed (\commit{d7d8b92}, \commit{8708c20}).}

\begin{reviewer}{Review notes}
Suggest adding User Manual section numbers to each step in the Description
field.
\end{reviewer}
\begin{response}
Done. Each tool's description now ends with ``(User's manual, Section~N.)''.
\end{response}
\status{Done (\commit{ae2e4d5}).}

% =====================================================================
\group{Documentation}

\begin{reviewer}{Review notes}
The documentation provided as a Word document does not match the
documentation linked from sedsat.org.
\end{reviewer}
\begin{response}
The Word document is the current manual; the sedsat.org copy is the earlier
version and will not be the published one. The manual source is now in the
repository, so there is one version.
\end{response}
\status{Done.}

\begin{reviewer}{Review notes}
``Panels'' are mentioned frequently. For novice users, explain what is meant by
panels; tell them to scroll down to see additional results.
\end{reviewer}
\begin{response}
A new section, ``The result window'' (Section~1.3.1), explains the panels, the
graph, table and export buttons, and scrolling for further panels.
\end{response}
\status{Done.}

\begin{reviewer}{Review notes}
More user guidance is needed in results interpretation, even if just references
to other work. A novice in statistics will have difficulty using this tool.
\end{reviewer}
\begin{response}
Added a section ``Interpreting the results'' (Section~4.3), with references to
your Sed\_SAT report and to the review by Collins and others (2017).
Interpretation is also added where the results are described: the
posterior predictive plots and what a measured value outside them means;
convergence of the MCMC samples; the DFA scatter plots; the error analysis
distributions; and source verification read as a classification check.
\end{response}
\status{Done.}

\begin{reviewer}{Review notes}
It would be helpful to include a table of method changes (or new methods) from
SedSat v1.0 to SedSat v3.0.
\end{reviewer}
\begin{response}
Added Table~1.1, ``Methods in SedSAT3'' (Section~1.2). For each method it
gives SedSAT3's approach, where the manual describes it, and how it compares
with Sed\_SAT~1.0: the same, changed, new, or not in SedSAT3. I wrote the
comparison from the Sed\_SAT~1.0 source code (the software release,
doi:10.5066/F76Q1VBX) and its report. Please check it, since you know
Sed\_SAT~1.0 best.

The comparison prompted two changes. SedSAT3 no longer corrects stable
isotopes for organic matter and size by default, as in Sed\_SAT~1.0, which
avoided biasing them (the user can still choose to correct them). Step~4 has a
confusion matrix (Comment~7.12). Sed\_SAT~1.0 also did three things that
SedSAT3 does not:
\begin{itemize}
    \item imputing values below detection;
    \item a tolerance in the bracket test;
    \item using tracer weighting factors from the DFA in the mixing model.
\end{itemize}
I would welcome your view on which of these SedSAT3 should have.
\end{response}
\status{Done; for your review.}

\begin{reviewer}{Review notes}
Figures are often too small to see without zooming in. Several naming
differences between the manual and the program.
\end{reviewer}
\begin{response}
Figures are now full page width (half width in side-by-side pairs). Most were
regenerated from the example data with the current program, often as the chart
or table alone rather than the whole screen. Every menu path and tool name in
the manual now matches the program exactly; the program's names were also
harmonized (``Multi-way'' throughout). A few screenshots of the main window
and of files still show the earlier version.
\mine{retake \code{SelectSourceTargetGroups.png}, \code{Constituent\_Viz.png},
\code{MCMC\_progress.png}, \code{MCMC\_folders.png}, \code{MCMC\_files.png},
\code{95p\_CI\_table.png} and \code{95p\_CI\_predicted.png} with the current
program.}
\end{response}
\status{Done, except the screenshots listed.}

% =====================================================================
\group{Comments in the user's manual: getting started and data}

\begin{reviewer}{Title page}
Version. Publication date. ``Missing'' (at the table of contents).
\end{reviewer}
\begin{response}
The title page now gives the software version, 1.1.8, and the date as a
placeholder until publication. The manual now also has a list of figures
after the table of contents. \todo{Cara: confirm the version number and set
the publication date.} I was not sure what ``missing'' at the table of
contents referred to: the list of figures, or a table of contents absent from
the Word copy (the PDF has one). Please tell me if it is something else.
\end{response}
\status{Done; version and date pending.}

\begin{reviewer}{Raw data preparation}
What about missing values or zeros? Are those allowed? How are negatives
managed? What if data is both positive and negative in a given element?
\end{reviewer}
\begin{response}
A new paragraph, ``Missing, zero and negative values'' (Chapter~2), explains:
\begin{itemize}
    \item Every cell must hold a finite number; an empty cell, text, NaN or Inf
          stops the import with the sheet, row and column. Values below
          detection must be replaced before import.
    \item Elements are modeled on a log scale, so an element with a zero or
          negative value cannot be used for fingerprinting; the program says
          which.
    \item Isotope $\delta$ values may be negative or of mixed sign: they are
          used on a linear scale and never log-transformed (Section~6.1.1,
          extended for this).
    \item The Box--Cox transformation leaves non-positive constituents
          untransformed and lists them.
\end{itemize}
\end{response}
\status{Done.}

\begin{reviewer}{Raw data preparation}
Tested by placing Aluminum in column 2 (passed); by changing Calcium to
Calcium\_µg/g and Arsenic to ``Arsenic space'' (passed). All edits conformed to
``each sheet must contain exactly the same elements''.
\end{reviewer}
\begin{response}
Thank you for these tests. The import also refuses a sheet whose header has a
gap, a duplicate sample name, and a blank row followed by more samples, each
with a message.
\end{response}

\begin{reviewer}{Raw data preparation}
\url{https://sedsat.org/wp-content/Data/SampleData.xlsx}: should this link be
updated to a code.usgs.gov address? What about virtual-mixture.xlsx?
\end{reviewer}
\begin{response}
The link is replaced by the example file's place in the repository and the
Windows package. On virtual-mixture.xlsx, see Section~3.
\end{response}
\status{Done.}

\begin{reviewer}{Importing the data}
Minimum number of samples per source group? Is there a minimum or maximum
number of source groups?
\end{reviewer}
\begin{response}
New paragraph ``Number of samples and source groups'' (Chapter~2):
\begin{itemize}
    \item At least two source groups, and at least two samples per group.
    \item The organic matter and size regression needs more samples in a group
          than coefficients.
    \item Discriminant function analysis needs more samples than constituents
          plus groups.
    \item The mass balance can separate at most one more source than the
          number of tracers.
\end{itemize}
There is no maximum number of groups other than this last limit.
\end{response}
\status{Done.}

\begin{reviewer}{Importing the data, Fig.~1.1}
My dialog box has blue circles and it's difficult to see what is selected and
what is not.
\end{reviewer}
\begin{response}
These are the Windows style's radio buttons; the selected one has a filled
centre. I agree that it is faint.
\mine{show the chosen target sheet in bold in that dialog.}
\end{response}
\status{To do.}

\begin{reviewer}{Importing the data}
The dialog box has an option to exclude elements, not samples. If ``exclude
all'' is clicked, the table does not update until the user clicks within the
table. Removing Lead was tricky: typing No gave a checked box that became Yes.
\end{reviewer}
\begin{response}
The manual now describes what happens: after import, the constituent properties
table opens, and samples are chosen separately in Settings → Include/Exclude
Samples. The ``Include all elements'' box now updates the table at once
(\commit{d7d8b92}). The \textit{Include in analysis} column is now a single
click on the box, with no editor to type into. That was fixed in the code
review work.
\end{response}
\status{Fixed.}

\begin{reviewer}{Importing the data}
It would be easier to read if the constituents appeared in the same order as in
the xlsx file.
\end{reviewer}
\begin{response}
Agreed. The program stores the constituents in a structure sorted by name, so
the alphabetical order appears in every table and result. Keeping the file's
order touches the project file format and every result. I would rather make
that change separately, with its own tests, than in this review. The manual now
states that constituents are listed alphabetically.
\mine{keep the Excel column order (project format and results); planned
separately.}
\end{response}
\status{Documented; change planned.}

\begin{reviewer}{Importing the data}
These columns can only be edited if the constituent type has been designated as
Isotope first; I suggest making this note in the manual. Please add a link to
further explanation in Mathematical Basis. Both N15 and C13 ratios for the
SampleData should be provided.
\end{reviewer}
\begin{response}
The manual now says to set the type first, and links to the stable isotope
section. It also says that the base element must be included in the analysis.
Both are now given in Section~2.1: N for N15 (ratio 0.003676) and TOC for C13
(ratio 0.011180, VPDB). TOC is the only carbon content measured, but it is also
the organic matter constituent, and the program did not accept that as a base
element. It now does. TOC is not used as a tracer. Fingerprinting uses it only
to weight each source's $\delta^{13}$C by its organic carbon content, with or
without the size and organic matter correction (\commit{67cc62f}). The
test project had used Calcium, which was not chemically meaningful. The
standard ratio cancels in the mixture's $\delta$ value, so its exact value does
not change the results.
\end{response}
\status{Done.}

\begin{reviewer}{Importing the data}
I don't see any values for C13; I assume the data is negative and the scale
does not allow for negative values.
\end{reviewer}
\begin{response}
Correct. The data plot always used a logarithmic axis. It now uses a linear
axis for a constituent with zero or negative values, and the manual says so.
\end{response}
\status{Fixed (\commit{d7d8b92}).}

\begin{reviewer}{Choosing samples to be discarded}
This did not automatically appear for me; I had to go to Settings and select
Include/Exclude Samples. Source groups appear in a drop-down menu, not tabs. I
do not see red highlight; worth noting that outliers are not automatically
removed. This statement should be moved up to where include/exclude is first
mentioned.
\end{reviewer}
\begin{response}
The section now matches the program:
\begin{itemize}
    \item the dialog is opened from Settings;
    \item groups are chosen from a drop-down list;
    \item a sample is excluded by clicking its Include cell, explained where
          the dialog is first described;
    \item flagged samples are not removed automatically.
\end{itemize}
On the red highlight: before the code review fixes, the notes that mark an
outlier were lost when the data were copied, so nothing could be red. With the
current code, a sample is red only when one of its values is beyond the
threshold. With the example data and a threshold of 3, no sample is, as the
regenerated figure shows; the manual now says so.
\end{response}
\status{Done.}

% =====================================================================
\group{Comments in the user's manual: deterministic fingerprinting}

\begin{reviewer}{Outlier analysis}
Box--Cox requires positive data: how is C13 handled? It has $\lambda=1$ for all
source groups. What if an element has both positive and negative values, as can
occur with N15?
\end{reviewer}
\begin{response}
A constituent with any zero or negative value is not transformed, and its
$\lambda$ is shown as 1, which is why C13 shows 1. The program lists the
constituents it left untransformed. For an element with mixed signs, such as
N15, the same applies to every group in which a non-positive value occurs.
This is now in Sections~3.1 and~6.1.1.
\end{response}
\status{Documented.}

\begin{reviewer}{Outlier analysis}
The execution window is confusing: it looks like you are running the outlier
test only on a single source group, but that is not the case. What is the
purpose of the source group drop-down?
\end{reviewer}
\begin{response}
There are two outlier tests:
\begin{itemize}
    \item The one run when you open Include/Exclude Samples covers every group.
    \item The tool in step~1 reports the deviations of the one group chosen in
          its form.
\end{itemize}
The section now explains both.
\end{response}
\status{Documented.}

\begin{reviewer}{Tools}
I'm finding the organization of the Tools window difficult to navigate. A user
new to fingerprinting will find this difficult. Consider ordering tools by
typical execution order rather than alphabetically. (Also: the order of the
descriptions should match the program; group the DFA tools; a Pre-Analysis
Tools subheader would be worthwhile.)
\end{reviewer}
\begin{response}
Done in the program and the manual.
\begin{itemize}
    \item The menu was alphabetical because the program's settings file format
          does not keep order. The menu now keeps the order written in the
          file.
    \item Other Statistical Tools is grouped in workflow order: Data
          exploration (outliers, corrected profiles, correlation, distribution
          fitting, K--S tests, Box--Cox); Element selection (ANOVA, auto-select,
          discriminant power); Discriminant function analysis (all six DFA
          tools together); Uncertainty and verification.
    \item A duplicate entry was removed.
    \item Chapter~5 of the manual follows the same groups and order, with the
          program's names.
\end{itemize}
\end{response}
\status{Done (\commit{d7d8b92}).}

\begin{reviewer}{Organic matter and particle size correction}
Program closed when the plot box was clicked. Are the corrected values plotted
anywhere? Some additional detail is needed: is size-corrected data used for the
organic correction, or raw data for both? How do you select the regression
option? What is ``Correct based on size and organic matter'' used for? Does
this need to be done for each source group? Is there a way to visualize raw vs
corrected data?
\end{reviewer}
\begin{response}
The crash is fixed (Section~1). The manual now explains, in Section~3.2:
\begin{itemize}
    \item When both are used, organic matter and particle size are the two
          variables of one multiple regression, fitted to the raw data; the
          correction applies both terms together, not one after the other.
    \item The regression form is chosen in the \textit{Equation} list
          (Linear or Power).
    \item One run of step~2 fits every source group.
    \item The regression plot shows the measured values.
    \item ``Correct based on size and organic matter'' fits nothing: it shows
          the corrected profiles for a chosen target sample, as a table or a
          chart. These can be compared with the raw values under Data, or
          exported.
\end{itemize}
\end{response}
\status{Documented; crash fixed.}

\begin{reviewer}{Organic matter and particle size correction}
Recommend expanding the axis so that points are clearly visible. Nothing is
highlighted in red for me (Figs.~2.4, 2.5). ``2a-Corrected results'' is not an
option; unable to check this section.
\end{reviewer}
\begin{response}
\begin{itemize}
    \item The regression plot's axes now extend beyond the data.
    \item The red p-values were lost by the same copy problem as the outlier
          flags; they show with the current code, as the regenerated figures
          show.
    \item The menu path is corrected: Other Statistical Tools → Data
          exploration → Correct based on size and organic matter.
\end{itemize}
\end{response}
\status{Fixed.}

\begin{reviewer}{Discriminant function analysis, step 4}
When selecting Box--Cox, Modify included elements, OM and Size for CTAIL1, I got
an error. This appears to be due to over-corrections: size and organic
corrections moved source values from positive to negative, which is an invalid
result. Only able to run DFA with power transform selected, or with no
corrections.
\end{reviewer}
\begin{response}
Your diagnosis is right. The error now names the element and group, and says
that the imported values are positive and the negative values come from the
correction. A note in Section~3.2 explains over-correction by the linear form
and recommends the power form, or not correcting the elements named.
\end{response}
\status{Message and documentation improved.}

\begin{reviewer}{Bracketing analysis}
To test, I raised CTAIL1/Aluminum to 10000; the test correctly failed. ``Value
higher'' — or lower? It would be nice to have a notes column for elements that
failed the bracket test. I do not understand the note about re-running the
bracket test for a new target sample. This does not look like bracket test
results (Fig.~2.7).
\end{reviewer}
\begin{response}
\begin{itemize}
    \item The text now reads ``higher than in all the source samples, or lower
          than in all of them''.
    \item The note is rewritten. The bracket test covers every target sample at
          once, but the constituent selection is one list for the project, so
          when you move to another target sample you adjust the list using
          that sample's column. For the batch tools, exclude every element that
          fails for any sample.
    \item The bracketing figure (now Fig.~3.7) was the wrong image; it is
          replaced by the bracket results table.
    \item A notes column is a good idea. Meanwhile, the manual suggests
          exporting the bracket table as the record. \mine{add a notes column
          to the constituent properties table.}
\end{itemize}
\end{response}
\status{Documented; notes column to do.}

\begin{reviewer}{Multi-way discriminant function analysis}
The benefits of multi-way DFA \ldots\ However, it is very sensitive to
multicollinearity and typically overfits if the number of variables $\gg$ the
number of samples. Given the structure of the typical sediment fingerprinting
input data, I think forward stepwise DFA would offer a better solution. Is data
tested/transformed for multivariate normality?
\end{reviewer}
\begin{response}
I agree, and step~4 is forward stepwise DFA. It adds one constituent at a time,
each the one that most improves the separation, and selects those up to the
smallest $\chi^2$ p-value. The tool's description said ``between the two
source groups'', which was wrong. It now describes forward stepwise selection
among all groups.

The manual now has a paragraph ``Stepwise selection'' (Section~3.4) giving
your reasons: correlated tracers, overfitting, and conditioning. On normality:
SedSAT3 does not test multivariate normality. Its Box--Cox option transforms
each constituent towards normality before the DFA, and the K--S tools check
constituents individually. Section~3.4 now says this.
\end{response}
\status{Documented.}

\begin{reviewer}{Multi-way discriminant function analysis}
Please explain the ``Modify the included elements based on the results'' option.
Is DFA run based on the size and organic corrections from only one target
sample? Each target sample has unique corrections and bracket results; each
could have unique DFA results.
\end{reviewer}
\begin{response}
Both are now explained in Section~3.5. With the correction, the DFA uses the
target sample chosen in the form, so the selection can differ between target
samples; for another target sample, run step~4 again with that sample.
\end{response}
\status{Documented.}

\begin{reviewer}{Multi-way discriminant function analysis}
Why is the x axis labeled ``Sample''? I did not understand the constituent order
until I looked at Wilks' lambda; perhaps put Wilks' lambda first. No y-axis label
on my DFA plots. My table gives p-values by element, not by source. Results in
Figs.~2.9--2.11 are found under Other Statistical Tools; if these are the most
relevant results they should appear in step~4.
\end{reviewer}
\begin{response}
\begin{itemize}
    \item The x axis now reads ``Constituents, in the order added'', and every
          plot has a y-axis title.
    \item The manual explains that all panels list the constituents in the
          order added. I kept the panel order, because the $\chi^2$ p-value is
          the selection criterion; Wilks' $\Lambda$ always decreases.
    \item You are right that the p-value table is per constituent. The
          per-group tables (Fig.~3.10) come from Discriminant function analysis
          (Multi-way), and the manual now says so.
    \item As you suggested, step~4 now also shows the discriminant scores and
          eigenvectors of the selected constituents (panels 5 and 6).
\end{itemize}
\end{response}
\status{Done (\commit{d7d8b92}).}

\begin{reviewer}{Multi-way discriminant function analysis, Fig.~2.10}
Figure needs updating, ellipses not present. What is the meaning of WB? Why do
the colors change with each pairing? Please verify color-blind accessibility per
USGS guidance. I see 2 WB\_Bank and no Pasture.
\end{reviewer}
\begin{response}
\begin{itemize}
    \item The colours were random, chosen again for each plot. All charts now
          use the Okabe--Ito colour-blind-safe palette, and in the DFA plot
          each group keeps its colour for every pair.
    \item ``WB\_'' was an internal prefix. The axes now read ``Discriminant
          score, Bank vs.\ the rest''.
    \item The window opened with the first group on both axes; it now starts
          with two different groups.
    \item Fig.~3.11 is regenerated, with the ellipses.
\end{itemize}
\todo{Cara: if USGS has its own colour guidance for figures, please point me to
it; I followed the Okabe--Ito palette.}
\end{response}
\status{Fixed (\commit{d7d8b92}).}

\begin{reviewer}{Multi-way discriminant function analysis, Fig.~2.11}
A confusion matrix would be a clearer output than this table. The goal is to
determine how well the tracers chosen can correctly assign sources, right?
\end{reviewer}
\begin{response}
Yes, and step~4 now has one: a seventh panel with the leave-one-out confusion
matrix of the source samples for the selected constituents. It uses linear
discriminant classification with equal priors, as in Sed\_SAT~1.0. Each row is
a true source group, each column the group assigned, and the last column the
percentage classified correctly. I checked it against an independent
calculation. In the example, most misclassifications are between Crop and
Pasture. The manual describes it and how to read it (Section~3.5).

Source Verification answers the same question with the fingerprinting itself,
and the manual now relates the two.
\end{response}
\status{Done.}

\begin{reviewer}{Fingerprinting}
Are the elements chosen by DFA the only ones used for fingerprinting, or are all
elements used? Is any weighting applied based on DFA results? Attempted to run:
program closed. Please add to the documentation that once the DFA step is run,
the constituent table will be modified.
\end{reviewer}
\begin{response}
Fingerprinting uses exactly the constituents marked as included, with equal
weight; the DFA weights are not used. Sections~3.5 and~3.6 now say so, with the
note about the table being modified. The crash is fixed (Section~1).
\end{response}
\status{Documented; crash fixed.}

\begin{reviewer}{Fingerprinting (Batch)}
If all samples receive the same corrections, which target sample should be used?
Specifically, what is used for $e_{i,t}$? Corrections are run for all target
samples, but DFA and Bayesian steps require selection of a target sample.
\end{reviewer}
\begin{response}
The batch tools correct for each target sample separately: $e_{i,t}$ is the
organic matter content and particle size of the target sample being analyzed.
What is shared is the constituent and sample selection, and that was made in
the earlier steps with one target sample. The note in Section~3.7 (repeated
for the Bayesian batch) now explains this. It suggests checking the selection
with target samples from both ends of the organic matter and size range.
\end{response}
\status{Documented.}

% =====================================================================
\group{Comments in the user's manual: Bayesian fingerprinting}

\begin{reviewer}{Bayesian Chemical Mass Balance Analysis}
8 chains is the default; is that generally recommended? 100,000 and 20,000
should be the defaults. The x-axis is labeled sample number, but the tick marks
are decimals. Using the current defaults the plot looks nothing like this:
samples appear random, no convergence.
\end{reviewer}
\begin{response}
\begin{itemize}
    \item The defaults are now 100,000 samples and 20,000 burn-in, so the
          default run converges. In the example, the chains settle by about
          15,000 to 20,000 samples (regenerated Fig.~4.2).
    \item The ticks are whole numbers.
    \item The manual explains the number of chains (one per processor core;
          more chains help judge convergence) and how to recognize convergence,
          and that too few samples look like scattered noise. The burn-in
          guidance is now consistent (20\%).
\end{itemize}
\end{response}
\status{Done (\commit{d7d8b92}).}

\begin{reviewer}{Bayesian Chemical Mass Balance Analysis}
The dashed lines from my run are far from the middle of the distribution, and in
one element not on the graph at all. Please give guidance on interpretation.
Add a figure for the fifth panel.
\end{reviewer}
\begin{response}
A measured value outside the plotted range is now always shown, the axis
extending to it. Section~4.1 now explains what a measured value in the tail
means:
\begin{itemize}
    \item a non-conservative tracer;
    \item a missing source;
    \item an unsuitable correction;
    \item a tracer that (nearly) failed the bracket test.
\end{itemize}
It also says what to do. Fig.~4.6 shows the fifth panel.
\end{response}
\status{Done.}

\begin{reviewer}{Bayesian Chemical Mass Balance Analysis (Batch)}
Add that the folders will be placed in the directory in which the project was
saved. I see 7 empty txt files; unable to verify content.
\end{reviewer}
\begin{response}
The manual now says where the folders go. They go to the working folder: the
folder of the project last opened or saved, or of the Excel file before the
project is saved. It advises saving the project first. With the current code
the seven files have content; I ran a batch on the example data and checked
each. The file names in the manual now match the program (\code{1\_MCMC
samples.txt}, \ldots).
\end{response}
\status{Fixed; screenshots of the folder to retake (Section~5).}

% =====================================================================
\group{Comments in the user's manual: other statistical tools}

\begin{reviewer}{Elements' correlation matrix}
Figure needs updating, current format below.
\end{reviewer}
\begin{response}
Regenerated. The result window no longer prints the whole matrix as text;
the table button shows it, with correlations beyond the threshold in red.
\end{response}
\status{Done.}

\begin{reviewer}{Analysis of Variance}
This assumes no corrections are made. Once OM and Size are corrected, Cobalt is
not listed as having low discriminating power; a discussion of this effect is
warranted. Red text not present in my table. D50 should not be part of this
analysis and should not be removed, possibly TOC as well.
\end{reviewer}
\begin{response}
\begin{itemize}
    \item I confirmed your observation: with the power correction for CTAIL1,
          Cobalt's p-value falls from 0.12 to about $10^{-9}$. A note now
          cautions against removing tracers without evaluating them on the data
          used for fingerprinting.
    \item The red text was the copy problem described above.
    \item D50 and TOC are no longer analysed by ANOVA, or by any tool that
          tests or selects tracers (\commit{d7d8b92}).
\end{itemize}
\end{response}
\status{Fixed and documented.}

\begin{reviewer}{Auto-select elements}
Colors repeated in graph: Bank and Crop identical to Forest and Pasture. D50
should not be part of this analysis; if the user checked Modify the included
elements, they would lose their size data.
\end{reviewer}
\begin{response}
The palette now has eight distinct colours, and each pair keeps its colour.
D50 and TOC are never tested or selected, and the \textit{Modify} option no
longer changes them. Your scenario (losing the size data) can no longer
happen.
\end{response}
\status{Fixed (\commit{d7d8b92}).}

\begin{reviewer}{Two-way DFA}
This is not on the list as ``Two-Way''; I assume this is one vs.\ one? x-axis
label is ``Sample''; same on other plots. F-test p value plot has no y axis label
or tick marks. Even when I uncheck ``Use only selected elements'' I only see 8
elements: running step 4 permanently altered my elements; there should be a note
in the manual and the program. It would be good to group these with the stepwise
DFA.
\end{reviewer}
\begin{response}
\begin{itemize}
    \item The sections now use the program's names: one vs.\ one, multi-way,
          and one vs.\ the rest. They are grouped with the stepwise versions,
          in the program and the manual (Section~5.3).
    \item The axis labels and the missing y axis are fixed (Section~4 of this
          letter).
    \item The note about step~4 is in the program and the manual.
    \item Unchecking ``Use only selected elements'' did not show the other
          elements because step~4 had already excluded them.
\end{itemize}
\end{response}
\status{Done.}

\begin{reviewer}{One vs.\ the rest DFA}
The weight vector: can this be explained more? Is it meant to be in the results?
\end{reviewer}
\begin{response}
Yes: the \textit{Eigen vector} panel shows it. Sections~3.4 and~5.3.3 now
explain that each sample's discriminant score is the weighted sum of its
constituent values, and how to read the weights.
\end{response}
\status{Documented.}

\begin{reviewer}{Distribution fitting}
This step did not run (all settings). Key shows observed and normal as the same
color; make the observed points black or another unused color.
\end{reviewer}
\begin{response}
The crash is fixed (Section~1). The samples are now black points, distinct from
both fitted curves.
\end{response}
\status{Fixed.}

\begin{reviewer}{Elements' Discriminant Power (Multi-way)}
In program ``multiple-way''. Program closed when running this step. Can
additional info be given as to why ``include target samples'' is an option?
\end{reviewer}
\begin{response}
\begin{itemize}
    \item The program now says ``Multi-way'', as the manual does.
    \item The crash is fixed (Section~1).
    \item The manual now explains the target option. It adds the target
          samples as another group, which shows the constituents that make the
          targets unlike a source. A constituent that separates the targets
          from every source is a warning, similar to a failed bracket test.
\end{itemize}
\end{response}
\status{Fixed and documented.}

\begin{reviewer}{Error Analysis}
Program closed (all settings). Why is there a black box and a green box in the
key?
\end{reviewer}
\begin{response}
\begin{itemize}
    \item The crash is fixed (Section~1).
    \item The two keys were the outline and the shaded area of the same
          distribution. The chart now has one key, named after the variable.
    \item The distribution panel now uses all realizations.
\end{itemize}
\end{response}
\status{Fixed.}

\begin{reviewer}{Kolmogorov--Smirnov tests}
These colors are hard to distinguish. When running lognormal for C13, observed
does not appear; I assume because of negatives? X axis label should be element,
constituent, or tracer, not sample.
\end{reviewer}
\begin{response}
\begin{itemize}
    \item The colours come from the new palette.
    \item A log-normal comparison of a constituent with non-positive values now
          stops with a message instead of drawing nothing.
    \item The axis reads ``Constituent''.
\end{itemize}
\end{response}
\status{Fixed (\commit{d7d8b92}, \commit{ae2e4d5}).}

\begin{reviewer}{Optimal Box-Cox parameters}
X axis label should be changed. A zero line for the y-axis would make this plot
easier to read.
\end{reviewer}
\begin{response}
Both done.
\end{response}
\status{Fixed.}

\begin{reviewer}{Source Verification}
Program closed (all options for Bank). Which OM/size correction will be used,
for which target sample?
\end{reviewer}
\begin{response}
The crash is fixed (Section~1). With the correction, each source sample being
tested is treated as the target, so the sources are corrected to that sample's
own organic matter content and particle size. The manual now says so.
\end{response}
\status{Fixed and documented.}

\begin{reviewer}{Genetic Algorithm estimation}
I do not see this option. I do not see any blue (modeled) in the modeled vs
observed plot.
\end{reviewer}
\begin{response}
The genetic algorithm is not part of this release: Levenberg--Marquardt solves
the same problem more efficiently, and the GA is only in an optional build. The
section now says so at its start. On the hidden modeled points: in all
modeled-vs-observed plots, the later series is now drawn smaller, so the
modeled value shows as a ring around the observed one where they agree.
\todo{Allen, Cara: keep the GA section (marked as not in this release), or
remove it from the manual?}
\end{response}
\status{Documented; plot fixed.}

\begin{reviewer}{Mathematical Basis}
Stable isotopes: how are negatives treated/resolved? N15 can span positive and
negative space. Some kind of ending or summary would be good; a header for
References is necessary.
\end{reviewer}
\begin{response}
Section~6.1.1 now explains how $\delta$ values of either sign are handled in
the mass balance and the likelihood. The chapter ends with a summary, and the
bibliography is titled ``References''.
\end{response}
\status{Done.}

% =====================================================================
\group{Items still open}

\begin{longtable}{@{}p{0.28\linewidth}p{0.18\linewidth}p{0.48\linewidth}@{}}
\toprule
Item & Who & Note \\
\midrule
\endhead
Repository path, DOI, version, publication date & \todo{Cara} & README, citation and manual title page hold placeholders. \\
sedsat.org and its download form & \todo{Cara, Allen} & Retire it or keep it as a pointer; check the form against USGS guidelines if it stays. Code-signing of the Windows package. \\
USGS colour guidance & \todo{Cara} & The charts use the Okabe--Ito colour-blind-safe palette. \\
Sed\_SAT 1.0 comparison & \todo{Allen, Lilly} & Review the last column of Table~1.1, and say which Sed\_SAT~1.0 features SedSAT3 should have. \\
Genetic algorithm section & \todo{Allen, Cara} & Keep (marked as not in this release) or remove. \\
Screenshots to retake & \mine{Arash} & Seven figures of the main window, the progress window and files. \\
virtual-mixture.xlsx & \mine{Arash} & Replace it with SampleData.xlsx in the Windows package. \\
Constituent order & \mine{Arash} & Keep the Excel column order; planned separately. \\
Notes column; target sheet dialog & \mine{Arash} & Enhancements suggested in the review. \\
\bottomrule
\end{longtable}

\end{document}
